\section*{Chapter \ref{ch:b1:stratigraphic-scale} --- The Stratigraphic Scale}

\begin{solution}{exo:b1:stratigraphic-scale:1}
(a)~Chronostratigraphic, a \omterm{def:b1:stratigraphic-scale:chronostratigraphy}{system}. (b)~Lithostratigraphic, a \omterm{def:b1:stratigraphic-scale:lithostratigraphy}{formation}.
(c)~Geochronologic, an epoch (its rock is the Lower Cretaceous \omterm{def:b1:stratigraphic-scale:chronostratigraphy}{Series}).
(d)~Biostratigraphic, a \omterm{def:b1:stratigraphic-scale:biozone}{biozone}. (e)~Chronostratigraphic, a \omterm{def:b1:stratigraphic-scale:chronostratigraphy}{stage}.
(f)~Geochronologic, an epoch. (g)~Lithostratigraphic, a \omterm{def:b1:stratigraphic-scale:lithostratigraphy}{member}.
\end{solution}

\begin{solution}{exo:b1:stratigraphic-scale:2}
Paleozoic $538.8 - 251.902 = \qty{286.9}{\Myr}$, \qty{53.2}{\%} of the
Phanerozoic; Mesozoic $251.902 - 66.00 = \qty{185.9}{\Myr}$, \qty{34.5}{\%};
Cenozoic \qty{66.0}{\Myr}, \qty{12.2}{\%}. The Phanerozoic is
$538.8/4567 = \qty{11.8}{\%}$ of the time since the base of the Hadean.
\end{solution}

\begin{solution}{exo:b1:stratigraphic-scale:3}
A, B, C, D were laid down in that order (superposition), horizontally (original
horizontality), so their tilt is a later event. The dyke E cuts A to D, so it is
younger than D (cross-cutting); whether it came before or after the tilting
cannot be told from the description. The erosion surface cuts D and E, so it is
younger than both (cross-cutting). F and G were then laid down on it
(superposition). The fault offsets A to F but not G: it moved after F and before G
(cross-cutting). Order: A, B, C, D; tilting and the intrusion of E (order
undecided); erosion; F; faulting; G.
\end{solution}

\begin{solution}{exo:b1:stratigraphic-scale:4}
$r = 30/(85.40 - 84.60) = \qty{37.5}{m}$ per million years. The first appearance
lies $12/37.5 = \qty{0.32}{\Myr}$ after the lower ash: $85.40 - 0.32 =
\qty{85.08}{\Ma}$.
\end{solution}

\begin{solution}{exo:b1:stratigraphic-scale:5}
(a)~Four vertical bars against height. (b)~B and C occur together from
\qtyrange{26}{50}{m}: a concurrent-range zone \qty{24}{m} thick. The \omterm{def:b1:stratigraphic-scale:biozone}{interval
zone} from the first appearance of B to that of D runs from \qtyrange{12}{41}{m},
\qty{29}{m}. (c)~In the second section the first appearances of B (\qty{3}{m})
and of C (\qty{9}{m}) can be placed, so the base of the \omterm{def:b1:stratigraphic-scale:biozone}{interval zone} of (b) and
the base of the concurrent-range zone are recognised; its top (the first
appearance of D) and the overlap's top as seen in section~1 cannot be checked,
D being absent. The first appearances of B and C ($\qty{12}{m} \leftrightarrow
\qty{3}{m}$, $\qty{26}{m} \leftrightarrow \qty{9}{m}$) are the safest lines; the
last appearances (B at \qty{50}{m} and \qty{15}{m}) depend more on preservation.
\end{solution}

\begin{solution}{exo:b1:stratigraphic-scale:6}
(a)~$60/20 = \qty{3}{\Myr}$: the top of the sandstone is three million years
younger in the landward well. (b)~One million years of mud at \qty{30}{m} per
million years: the ash lies \qty{30}{m} above the top of the sandstone.
(c)~When the ash fell, the shoreline had moved \qty{20}{km} inland of the seaward
well, so it was still \qty{40}{km} seaward of the landward well, which stood on
dry land. The ash fell on the land surface there, below the sandstone that the
sea laid down two million years later --- if it was not washed away first, as a
loose layer on land may well be. The same ash bed is above the \omterm{def:b1:stratigraphic-scale:lithostratigraphy}{formation} boundary in
one well and below it in the other: the boundary crosses time.
\end{solution}

\begin{solution}{exo:b1:stratigraphic-scale:7}
(a)~Lochkovian $1.36/419.62 = \qty{0.32}{\%}$; Induan $0.024/251.902 =
\qty{0.0095}{\%}$, about 34 times smaller. The base of the Induan lies among ash
beds dated precisely in the GSSP section itself; a large uncertainty suggests a
boundary whose age must be carried from dated levels farther away.
(b)~$251.941 - 251.880 = \qty{0.061}{\Myr} = 61$ thousand years, with the
uncertainty $\sqrt{37^2 + 31^2} = 48$ thousand years: the published
\num[uncertainty-mode=separate]{60 \pm 48} thousand years.
\end{solution}

\begin{solution}{exo:b1:stratigraphic-scale:8}
(a)~$m = 3000 \times \frac43\pi\,(5000)^3 = \qty{1.57e15}{kg}$; the dust,
$60m = \qty{9.42e16}{kg}$. (b)~The Earth's surface is $4\pi(6.378\times10^6)^2
= \qty{5.11e14}{m^2}$: \qty{184}{kg} per square metre. (c)~$184/2500 =
\qty{0.074}{m}$, about \qty{7}{cm}. (d)~The thickness goes as the cube of the
diameter: \qty{1.6}{cm} for \qty{6}{km}, \qty{20}{cm} for \qty{14}{km}. The order
of magnitude, centimetres to decimetres, is that of the boundary clay,
\qty{50}{cm} thick at El Kef; the comparison can only be rough, since the dust
cannot have fallen evenly and a clay also gathers the ordinary sediment of its
time.
\end{solution}

\begin{solution}{exo:b1:stratigraphic-scale:9}
(a)~$\num{60e3} \times \num{2e-6} = \qty{0.12}{m}$, \qty{12}{cm}; for 12 and 108
thousand years, \qty{2.4}{cm} and \qty{21.6}{cm}. (b)~$\num{32e3} \times
\num{10e-6} = \qty{0.32}{m}$. (c)~An extinction occupies centimetres to decimetres
of rock: samples must be taken centimetre by centimetre, and a single sample
\qty{10}{cm} thick in the condensed section averages fifty thousand years ---
most of the event.
\end{solution}

\begin{solution}{exo:b1:stratigraphic-scale:10}
(a)~$u = R/s = 10 \times 2000 = \qty{20000}{m}$ per million years,
\qty{20}{km} per million years; over \qty{30}{km}, $\Delta t = \qty{1.5}{\Myr}$.
(b)~The plain is now the surface $z = s x + a t$. The shoreline satisfies
$s x + a t = z_0 + R t$, so $x(t) = (z_0 + (R - a)t)/s$ and
$\mathrm dx/\mathrm dt = (R - a)/s$, positive (inland) while $a < R$. When
$a > R$ the speed is negative: the shoreline moves seaward, a regression, and
the shoreline sandstone is then younger seaward. (c)~A cliff \qty{200}{m} long spans,
at \qty{20}{km} per million years, $0.2/20 = \qty{0.01}{\Myr}$ --- ten thousand
years of diachronism, far below what fossils or most dates can resolve; the rock on
either side of the contact is the same everywhere along it.
\end{solution}

\begin{solution}{exo:b1:stratigraphic-scale:11}
(a)~$f = (25 - 10)/(50 - 10) = 0.375$ and $t = 100.50 - 0.375 \times 0.80 =
\qty{100.20}{\Ma}$, with
\[
  \sigma_t = \sqrt{(0.625 \times 0.05)^2 + (0.375 \times 0.03)^2} = \qty{0.033}{\Myr}.
\] (b)~$g(f) = (1 - f)^2\sigma_1^2 + f^2\sigma_2^2$
has $g'(f) = -2(1 - f)\sigma_1^2 + 2f\sigma_2^2$, zero at $f^* =
\sigma_1^2/(\sigma_1^2 + \sigma_2^2) = 0.0025/0.0034 = 0.735$, nearer the better
dated bed; there $\sigma_t = \sigma_1\sigma_2/\sqrt{\sigma_1^2 + \sigma_2^2} =
\qty{0.026}{\Myr}$, smaller than either. (c)~$g$ is a parabola opening upward, so
on $[0, 1]$ it is largest at an end: $g(0) = \sigma_1^2$, $g(1) = \sigma_2^2$.
Hence $\sigma_t \le \max(\sigma_1, \sigma_2)$ between the beds.
\end{solution}

\begin{solution}{exo:b1:stratigraphic-scale:12}
(a)~The boundary is halfway in depth: $f = 0.5$, \qty{29.50}{\Ma}. (b)~With
$\varphi_0\lambda = 0.63 \times 1960 = \qty{1234.8}{m}$, the solid thickness from
\qtyrange{600}{1000}{m} is $400 - 1234.8\,(e^{-600/1960} - e^{-1000/1960}) =
\qty{232.2}{m}$, and from \qtyrange{200}{600}{m}, \qty{194.2}{m}; in all
\qty{426.3}{m}. The fraction of the grains laid down between the lower ash and the
boundary is $232.2/426.3 = 0.545$, and the age is $30.0 - 0.545 =
\qty{29.46}{\Ma}$. (c)~The second estimate is \num{45} thousand years younger.
It is the better one: the deeper half has lost more of its water, so it holds
more grains, and more time, per metre than the shallower half (mean porosity
\num{0.42} against \num{0.51}).
\end{solution}

\subsection*{Problem: Correlating three sections across a basin}

\begin{solution}{pb:b1:stratigraphic-scale:1}
\begin{enumerate}
  \item Three columns at, say, \qty{1}{cm} for \qty{10}{m}, joined by the ash
        beds and by the zone boundaries; the Kell Sandstone wedge thins from
        \qty{45}{m} in A to \qty{6}{m} in C.
  \item Zone~2: A $65.8 - 36.6 = \qty{29.2}{m}$, B \qty{20.8}{m}, C
        \qty{10.0}{m}. Zone~3 logged: A \qty{14.2}{m}, B \qty{14.0}{m}, C
        \qty{7.2}{m}.
  \item Zone~1 begins at the first appearance of \emph{V.~primum}, which is found
        from the base of every section: its first appearance lies below them.
  \item It lives only in the sandstone: it follows the facies, and the facies
        is of different ages from place to place. In A it occurs in zones~1
        and~2, in B and C only in zone~1.
  \item In zone~2 in all three (A \qty{49.0}{m}, B \qty{40.0}{m}, C
        \qty{19.0}{m}). It is consistent from section to section here, and so
        could subdivide zone~2, but a last appearance depends on preservation;
        it should be checked in more sections first.
  \item The lines join \qty{36.6}{m}, \qty{31.2}{m} and \qty{14.8}{m}, then
        \qty{65.8}{m}, \qty{52.0}{m} and \qty{24.8}{m}. Zone~2 thins in the
        ratio $29.2 : 20.8 : 10.0$, about $2.9 : 2.1 : 1$: if the zone lasted
        the same time everywhere, the sediment accumulated about three times
        faster in A than in C.
  \item A: zone~2 ($45.0 > 36.6$); B: zone~1 ($20.0 < 31.2$); C: zone~1
        ($6.0 < 14.8$).
  \item No: it is in zone~2 in A, zone~1 in B and C, so it is younger in the
        west. The marine shale spread westward over the shoreline sand: the
        shoreline moved west, inland, in a transgression.
  \item An ash falls over the whole basin in a geological instant and settles
        on any sediment; the top of the sandstone follows the moving shoreline
        and is \omterm{def:b1:stratigraphic-scale:diachronous}{diachronous}.
  \item Lateral continuity (the bed continues across the basin), and
        superposition (ash~1 below ash~2 in every section, which checks the
        correlation).
  \item A $(71.0 - 5.0)/1.20 = \qty{55.0}{m}$ per million years; B
        $48.0/1.20 = \qty{40.0}{}$; C $24.0/1.20 = \qty{20.0}{}$.
  \item A is nearest the land, where the sediment was supplied; less of it
        reached the deeper water to the east.
  \item $0.1/55 = \num{1.8}$ thousand years in A, $\num{2.5}$ in B, $\num{5.0}$
        in C: A resolves time most finely.
  \item A $152.40 - 40.0/55.0 = \qty{151.67}{\Ma}$; B $152.40 - 12.0/40.0 =
        \qty{152.10}{\Ma}$; C $152.40 - 3.0/20.0 = \qty{152.25}{\Ma}$.
  \item C to B: $20/(152.25 - 152.10) = \qty{133}{km}$ per million years; B to A:
        $25/(152.10 - 151.67) = \qty{58.5}{km}$ per million years. It slowed
        down.
  \item A $152.40 - 60.8/55.0 = \qty{151.29}{\Ma}$; B $152.40 - 44.0/40.0 =
        \qty{151.30}{\Ma}$; C $152.40 - 21.8/20.0 = \qty{151.31}{\Ma}$.
  \item They span \num{16} thousand years, less than the uncertainty of each
        (about \num{28} thousand years, question~19): the first appearance is
        synchronous across the basin within the resolution of the dates, which
        is what makes it a good index.
  \item A $152.40 - 31.6/55.0 = \qty{151.826}{\Ma}$; B $152.40 - 23.2/40.0 =
        \qty{151.820}{\Ma}$; C $152.40 - 11.8/20.0 = \qty{151.810}{\Ma}$.
  \item $f = 0.479$, $0.483$, $0.492$; $\sigma_t = \sqrt{(1 - f)^2 + f^2} \times
        0.04 = \qty{0.028}{\Myr}$ in each, \num{28} thousand years.
  \item Mean \qty{151.819}{\Ma}; spread $151.826 - 151.810 = \num{16}$ thousand
        years, smaller than the uncertainty of each age: the three agree.
  \item The interpolation assumes one rate from ash~1 to ash~2. In A that
        interval holds shoreline sandstone up to \qty{45}{m} and shale above;
        sand and mud need not have accumulated at the same rate, and they
        compact differently (\cref{ch:b1:diagenesis-basins}). If the sand
        accumulated faster, the sandstone part holds less time than its
        thickness suggests, and the first appearance is older than computed
        --- the age from A is indeed the oldest of the three.
  \item B: the first appearance lies in marine shale, \qty{11}{m} above the
        top of the sandstone, with no change of facies near the level, and the
        rate gives \num{2.5} thousand years per \qty{10}{cm}. A resolves time
        better but places the level in the shoreline sandstone, \qty{8.4}{m}
        below a change of facies; C is condensed, \num{5} thousand years per
        \qty{10}{cm}. C is the worst for resolution, A for facies.
  \item From the mean ages of the two first appearances, $151.819 - 151.302 =
        \qty{0.52}{\Myr}$: about half a million years.
  \item The boundary is \qty[uncertainty-mode=separate]{151.82(3)}{\Ma}.
\end{enumerate}
\end{solution}
